% ── §17 Discussion ──
\section{Discussion}
\label{sec:discussion}

\subsection{Summary of Contributions}
This work establishes that the analytical chain from biological sequences
through Gray-code encoding, Karnaugh-map construction, Quine--McCluskey
minimisation, and prime-implicant extraction can be made machine-provable in
Lean~4, and that the resulting Boolean circuits carry genuine structural
information when applied to protein multiple sequence alignments.  The key
findings are: (1) universal contact circuits transfer across proteins at
$5.1\times$ random enrichment; (2) Gray-code adjacency is enriched among
native contacts with $z = 2.51$ specificity for the particular encoding;
(3) the Karnaugh-map cell-rate prior improves global DCA ranking and, when
fused via rank aggregation, exceeds both individual methods; (4) flipped
(forbidden-combination) circuits transfer with 98.3\% specificity; and (5)
conserved column pairs are enriched for contacts by $3.2\times$, confirming
that the conserved structural core carries information through constraint
rather than through covariation.

\subsection{Comparison with Literature}
Our mfDCA implementation achieves prec@L/5 = 0.169 on PSICOV150, consistent
with published mean-field benchmarks of approximately 0.15--0.20 on this
dataset.  The substantial gap to PSICOV's sparse inverse covariance method
(0.727) reflects the fundamental difference between simple matrix inversion
and $\ell_1$-regularised estimation, not an implementation error.  Our ranksum
fusion (0.204) demonstrates that the K-map prior contributes orthogonal
information to any continuous scoring method, though it does not close this
methodological gap.

The finding that Gray adjacency is enriched among native contacts ($1.178\times$,
$p < 10^{-77}$) resolves an ambiguity from earlier underpowered analysis on a
single protein family.  The permutation control ($z = 2.51$ against random
code assignments) confirms that the specific He/Gray structure carries genuine
signal beyond what generic physicochemical clustering would produce.

\subsection{Limitations}
Several limitations constrain the scope of our conclusions.  First, consensus
residue features collapse within-column variation into a single code per
position, discarding frequency information that richer encodings might
exploit.  Second, separation bins are globally defined without per-protein
length normalisation.  Third, the reconstruction assessment is limited to six
proteins by the computational cost of exact Quine--McCluskey on large padded
maps.  Fourth, the cross-dataset generalisation experiment on EVmutation was
inconclusive due to A2M format parsing complexity, leaving open the question
of whether findings replicate beyond PSICOV150.

\subsection{Future Directions}
Three directions emerge naturally from this work.  First, applying the complete
circuit framework to CASP free-modelling targets with published contact maps
would test whether the rank-fusion improvement holds on harder targets.
Second, replacing binary conservation flags with continuous entropy bins as
additional Karnaugh-map dimensions could capture finer-grained structural
information about the buried core.  Third, integrating circuit-derived rules
as input features to deep learning models (rather than using them standalone)
might combine the interpretability of Boolean logic with the representational
power of neural networks.

\subsection{Data and Code Availability}
All source code, formal proofs, result files, and documentation are available
at \url{https://github.com/E-Motioner-X-SBS/co-evolution-analysis}.  The Lean~4
formalisation is available at \url{https://github.com/E-Motioner-X-SBS/lean_proofs}.
The PSICOV150 dataset is available from \url{https://bioinfadmin.cs.ucl.ac.uk/downloads/PSICOV/}.
